\emph{Preliminaries.} Write $\tau = \langle\cdot\rangle_{0}$ and $\mathcal{A} = \Cl_{k}(p,q)$. By Theorem \claimref{synthesis.signed-directions-and-the-integer-clifford-algebra} the blades are a basis, each blade is invertible with inverse $\pm$ itself, and two blades commute or anticommute. The blade inverse $\dagger$ is an anti-automorphism with $\tau(x^{\dagger}) = \tau(x)$, $\tau(xy) = \tau(yx)$, and $\tau(x^{\dagger}x)$ is the sum of the squares of the coefficients of $x$ (clause (iii) of Theorem \claimref{synthesis.chsh-bound-for-two-subsystems-of-a-clifford-algebra-by-a-rational-certificate}), so $\tau(x^{\dagger}x) > 0$ for $x \ne 0$. A blade $b$ is self-adjoint exactly when $b^{2} = 1$ and anti-self-adjoint exactly when $b^{2} = -1$, so the self-adjoint elements are the span of the blades of square $1$. The corner of a projector $\Pi$ is closed under $\dagger$, so each of its elements is the sum of a self-adjoint and an anti-self-adjoint element of the corner; if $\Pi$ is an atom, every element of the corner is $c\Pi + a$ with $c \in k$ and $a^{\dagger} = -a$, and then $\tau(a) = \tau(a^{\dagger}) = -\tau(a)$ gives $\tau(a) = 0$.

\emph{(i).} Let $G_{b}$ be the set of products of subsets of $b_{1}, \dots, b_{m}$; by independence its $2^{m}$ elements are $\pm$ distinct blades and only the empty product is $\pm 1$. Since the $b_{i}$ commute and square to $1$, the factors $\tfrac12(1 + \chi_{i}b_{i})$ are commuting projectors, so $\Pi_{\chi}$ is a projector; expanding, $\Pi_{\chi} = 2^{-m}\sum_{g \in G_{b}}\pm g$, whose scalar part is $2^{-m}$, so $\Pi_{\chi} \ne 0$. For $\chi \ne \chi'$ the product $\Pi_{\chi}\Pi_{\chi'}$ contains a factor $\tfrac12(1 + b_{i})\tfrac12(1 - b_{i}) = 0$, and $\sum_{\chi}\Pi_{\chi} = \prod_{i}\bigl(\tfrac12(1 + b_{i}) + \tfrac12(1 - b_{i})\bigr) = 1$. The corner: let $b$ be a blade. If $b$ anticommutes with some $b_{i}$, then $\tfrac12(1 + \chi_{i}b_{i})\,b\,\tfrac12(1 + \chi_{i}b_{i}) = b\,\tfrac12(1 - \chi_{i}b_{i})\tfrac12(1 + \chi_{i}b_{i}) = 0$, so $\Pi_{\chi}b\Pi_{\chi} = 0$. Otherwise $b$ commutes with every $b_{i}$ and $\Pi_{\chi}b\Pi_{\chi} = b\Pi_{\chi}$. So the corner is spanned by the $b\Pi_{\chi}$ with $b$ commuting with all $b_{i}$. For such $b$: $(b\Pi_{\chi})^{\dagger} = \Pi_{\chi}b^{-1} = b^{-1}\Pi_{\chi}$. If $b^{2} = 1$, maximality gives $b = \pm g$ with $g \in G_{b}$, and $g\Pi_{\chi} = \pm\Pi_{\chi}$ because $\chi_{i}b_{i}\tfrac12(1 + \chi_{i}b_{i}) = \tfrac12(1 + \chi_{i}b_{i})$; so $b\Pi_{\chi} = \pm\Pi_{\chi}$. If $b^{2} = -1$, then $b^{-1} = -b$, so $b\Pi_{\chi}$ is anti-self-adjoint, and $(b\Pi_{\chi})^{2} = b^{2}\Pi_{\chi} = -\Pi_{\chi}$. Hence the corner is $k\Pi_{\chi}$ plus a span of anti-self-adjoint elements, and its self-adjoint part is $k\Pi_{\chi}$: $\Pi_{\chi}$ is an atom. The conditions of independence and maximality do not change when signs are changed. Existence: the empty set is independent, the number of independent commuting blades is bounded, and an independent commuting set that is not maximal can be enlarged by a blade $b$ of square $1$ commuting with it and not in $\pm G_{b}$, the enlarged set being independent because a product equal to $\pm 1$ would put $b$ in $\pm G_{b}$. Modes: a blade commutes with $p_{j} = f_{j}e_{j}$ exactly when it contains both or neither of $e_{j}$, $f_{j}$, so a blade commuting with every $p_{j}$ is $\pm$ a product of parities, the parities are independent (their products are distinct blades), and the corner of $\Pi_{0}$ is $k\Pi_{0}$. With $v_{j} = \tfrac12(e_{j} + f_{j})$ of Proposition \claimref{tower.ladder-pairs-satisfy-car-vectors}: $v_{j}p_{j} = \tfrac12(e_{j}f_{j}e_{j} + f_{j}f_{j}e_{j}) = \tfrac12(-f_{j} - e_{j}) = -v_{j}$, so $v_{j}(1 + p_{j}) = 0$, and $v_{j}$ commutes with $p_{i}$ for $i \ne j$ (an even blade on other generators), so $v_{j}\Pi_{0} = 0$. The table of dimensions is the script's.

\emph{(ii).} $\Pi y^{\dagger}y\Pi$ is a self-adjoint element of the corner of the atom $\Pi$, so it is $n_{y}\Pi$; taking scalar parts, $n_{y}\tau(\Pi) = \tau(\Pi y^{\dagger}y\Pi) = \tau((y\Pi)^{\dagger}(y\Pi)) > 0$ when $y\Pi \ne 0$, and $\tau(\Pi y^{\dagger}y\Pi) = \tau(\Pi y^{\dagger}y)$. Put $R = y\Pi y^{\dagger}/n_{y}$. Then $R^{\dagger} = R$, $R^{2} = y(\Pi y^{\dagger}y\Pi)y^{\dagger}/n_{y}^{2} = R$, and $\tau(R) = \tau(\Pi y^{\dagger}y)/n_{y} = \tau(\Pi)$. If $r$ is a self-adjoint element of the corner of $R$, then $r = RrR = y\,\Pi(y^{\dagger}ry)\Pi\,y^{\dagger}/n_{y}^{2}$, and $\Pi(y^{\dagger}ry)\Pi$ is self-adjoint in the corner of $\Pi$, hence a multiple of $\Pi$, so $r$ is a multiple of $R$: $R$ is an atom. If the corner of $\Pi$ is $k\Pi$, the same computation with any $x$ in place of $r$ gives $RxR \in kR$, so the corner of $R$ is $kR$, and $RxR = cR$ with $c\tau(R) = \tau(RxR) = \tau(Rx)$. For the decomposition, $yy^{\dagger} = y\bigl(\sum_{\chi}\Pi_{\chi}\bigr)y^{\dagger} = \sum_{\chi}y\Pi_{\chi}y^{\dagger}$, and each term is zero (when $y\Pi_{\chi} = 0$) or $n_{y,\chi}(\Pi_{\chi})_{y}$. A linear functional non-negative on atoms is therefore non-negative on $yy^{\dagger}$; conversely an atom $\Pi$ is $\Pi\Pi^{\dagger}$.

\emph{(iii).} Here $\omega$ is a state in the sense of Premise \claimref{synthesis.premise-positivity-of-values}; the three clauses are used where named. Let $N = 1 - \Pi$, a projector, so $N = NN^{\dagger}$ and $\omega(NN^{\dagger}) = 1 - \omega(\Pi) = 0$. The form $\beta(x,y) = \omega(xy^{\dagger})$ is bilinear, symmetric by the reality clause ($\omega(yx^{\dagger}) = \omega((xy^{\dagger})^{\dagger}) = \omega(xy^{\dagger})$) and $\beta(x,x) \ge 0$ by positivity. If $\beta(N,y) \ne 0$ for some $y$, take $t = -\beta(N,y)/(\beta(y,y) + 1) \in k$: then $\beta(N + ty, N + ty) = 2t\beta(N,y) + t^{2}\beta(y,y) = -\beta(N,y)^{2}\,(\beta(y,y) + 2)/(\beta(y,y) + 1)^{2} < 0$, a contradiction; this uses only that $k$ is ordered. So $\omega(Ny^{\dagger}) = 0$ for all $y$, that is $\omega(Nx) = 0$ for all $x$, and $\omega(xN) = \omega((xN)^{\dagger}) = \omega(Nx^{\dagger}) = 0$. Therefore $\omega(x) = \omega((\Pi + N)x(\Pi + N)) = \omega(\Pi x\Pi)$. Write $\Pi x\Pi = c\Pi + a$ with $a^{\dagger} = -a$ (Preliminaries); reality gives $\omega(a) = -\omega(a) = 0$, so $\omega(x) = c$, and $\tau(\Pi x) = \tau(\Pi x\Pi) = c\tau(\Pi)$. For a projector $E$: $\tau(\Pi E) = \tau(\Pi EE\Pi) = \tau((E\Pi)^{\dagger}(E\Pi)) \ge 0$, and $\tau(\Pi) = \tau(\Pi E) + \tau(\Pi(1 - E))$ with both terms non-negative, so the weight is in $[0,1]$; additivity over a family with sum $1$ is linearity. For the orbit: $\tau(\Pi_{x}\Pi_{y}) = \tau(x\Pi x^{\dagger}y\Pi y^{\dagger})/(n_{x}n_{y}) = \tau(\Pi x^{\dagger}y\Pi\,\Pi y^{\dagger}x\Pi)/(n_{x}n_{y}) = \tau(zz^{\dagger})/(n_{x}n_{y})$, and dividing by $\tau(\Pi_{x}) = \tau(\Pi)$ gives the formula, symmetric in $x$, $y$ because $\tau(zz^{\dagger}) = \tau(z^{\dagger}z)$. If the corner is $k\Pi$ then $z = c\Pi$ and $\tau(zz^{\dagger}) = c^{2}\tau(\Pi)$. Since $\Pi_{y}$ is unchanged when $y$ is multiplied by a nonzero element of $k$, no normalisation of $y$ is needed.

\emph{(iv).} (a) Positivity and reality of $\omega_{\Pi}$ are clause (iii) of the theorem cited. The functional $x \mapsto \tau(f_{1}x)$ vanishes on self-adjoint $x$: $\tau(f_{1}x) = \tau((f_{1}x)^{\dagger}) = \tau(x^{\dagger}f_{1}^{\dagger}) = -\tau(xf_{1}) = -\tau(f_{1}x)$. Every $xx^{\dagger}$, $1$ and $\Pi$ are self-adjoint, so $\mu$ agrees with $\omega_{\Pi}$ on them; $\mu(f_{1}) = \tau(f_{1}f_{1}) = -1$. That in general the conclusion holds on self-adjoint elements without reality: apply the argument of (iii) to the symmetric part $\tfrac12(\beta(x,y) + \beta(y,x))$, which gives $\omega(Ny^{\dagger} + yN) = 0$ for all $y$; for self-adjoint $x$ put $y = \Pi x$ and $y = \tfrac12 NxN$ to get $\omega(Nx\Pi + \Pi xN) = 0$ and $\omega(NxN) = 0$, so $\omega(x) = \omega(\Pi x\Pi) = c$, since $\Pi x\Pi$ is then self-adjoint in the corner. (b) $1 + p_{1} + \tfrac13 e_{1}$ is self-adjoint, so $\nu$ is real; the values are the script's. (c) Let $s$ be self-adjoint in the corner of $\Pi'$ with $\tau(s) = 0$, $s \ne 0$ (if the self-adjoint part of the corner is larger than $k\Pi'$, subtract the multiple of $\Pi'$ with the same scalar part). Put $\varphi_{t}(x) = \tau((\Pi' + ts)x)/\tau(\Pi')$. It is real because $\Pi' + ts$ is self-adjoint; $\varphi_{t}(1) = 1 + t\tau(s)/\tau(\Pi') = 1$ and $\varphi_{t}(\Pi') = \tau(\Pi' + ts\Pi')/\tau(\Pi') = 1$, as $s\Pi' = s$. Positivity: $\tau(\Pi')\varphi_{t}(yy^{\dagger}) = \tau(y^{\dagger}(\Pi' + ts)y) = q_{0}(u) + tq_{1}(u)$ with $u = \Pi'y$, $q_{0}(u) = \tau(u^{\dagger}u)$ and $q_{1}(u) = \tau(u^{\dagger}su)$ (using $s = \Pi's\Pi'$). On the space $U = \Pi'\mathcal{A}$ the form $q_{0}$ is positive definite, so $U$ has a basis $u_{1}, \dots, u_{r}$ orthogonal for it (the Gram-Schmidt process needs only division by the nonzero numbers $q_{0}(u_{i}) = d_{i} > 0$). With $u = \sum_{i}c_{i}u_{i}$ and $M$ the largest absolute value of an entry of the matrix of $q_{1}$ in this basis, $|q_{1}(u)| \le M\sum_{i,j}|c_{i}||c_{j}| \le Mr\sum_{i}c_{i}^{2} \le \varepsilon^{-1}q_{0}(u)$ with $\varepsilon = \min_{i}d_{i}/(Mr)$ (any $\varepsilon > 0$ if $M = 0$). So $\varphi_{t}(yy^{\dagger}) \ge 0$ for $|t| \le \varepsilon$: $\varphi_{t}$ is a state sharp on $\Pi'$. It differs from $\varphi_{0}$ at $x = s$: $\varphi_{t}(s) = t\tau(s^{2})/\tau(\Pi')$, nonzero for $t \ne 0$ since $\tau(s^{2}) = \tau(s^{\dagger}s) > 0$, and $\varphi_{0}(s) = \tau(s)/\tau(\Pi') = 0$. Conversely if $\Pi'$ is an atom the state is unique by (iii). In the example $s = \tfrac12(1 + p_{1})p_{2}$ has $\tau(s) = 0$, and $\Pi' + ts = \tfrac{1 + t}{2}\,\Pi'(1 + p_{2}) + \tfrac{1 - t}{2}\,\Pi'(1 - p_{2})$ is for $|t| \le 1$ a combination with non-negative coefficients of the elements $\Pi'(1 \pm p_{2})$, each twice a projector; $\varphi_{t}(p_{2}) = t$.
